\section{Subtree frequency queries}

DSU on Tree (also called \emph{Sack}) efficiently maintains information for every subtree. A classic task is: for every node \texttt{u}, find the maximum frequency of a colour in its subtree and the sum of all colours having that frequency. For example, \texttt{[1, 2, 2, 3, 3, 3]} has maximum frequency 3 and answer 3; \texttt{[1, 1, 2, 2, 3]} has maximum frequency 2 and answer $1+2=3$.

The typical complexity is $O(n \log n)$ time and $O(n)$ space. The key is to retain data from the largest child subtree while discarding and rebuilding smaller subtrees.

\subsection{Most frequent colour in every subtree}
\begin{lstlisting}
int n;
vector<int> color, sz, heavy, cnt;
vector<vector<int>> g;
vector<long long> ans;
int mxFreq = 0;
long long sumOfColors = 0;

void dfs_sz(int u, int p) {
    sz[u] = 1; heavy[u] = -1;
    int bestSize = 0;
    for (int v : g[u]) if (v != p) {
        dfs_sz(v, u);
        sz[u] += sz[v];
        if (sz[v] > bestSize) bestSize = sz[v], heavy[u] = v;
    }
}

void addColor(int c) {
    cnt[c]++;
    if (cnt[c] > mxFreq) mxFreq = cnt[c], sumOfColors = c;
    else if (cnt[c] == mxFreq) sumOfColors += c;
}

void addSubtree(int u, int p) {
    addColor(color[u]);
    for (int v : g[u]) if (v != p) addSubtree(v, u);
}

void clearSubtree(int u, int p) {
    cnt[color[u]]--;
    for (int v : g[u]) if (v != p) clearSubtree(v, u);
}

void dfs_dsu(int u, int p, bool keep) {
    // Light children: compute their answers, then discard their data.
    for (int v : g[u])
        if (v != p && v != heavy[u]) dfs_dsu(v, u, false);

    // Heavy child: keep its data for reuse.
    if (heavy[u] != -1) dfs_dsu(heavy[u], u, true);

    // Add all light subtrees to the heavy child's existing data.
    for (int v : g[u])
        if (v != p && v != heavy[u]) addSubtree(v, u);

    addColor(color[u]);          // data now represents subtree(u)
    ans[u] = sumOfColors;

    if (!keep) {
        clearSubtree(u, p);
        mxFreq = 0;
        sumOfColors = 0;
    }
}

int main() {
    ios::sync_with_stdio(false);
    cin.tie(nullptr);

    cin >> n;
    color.resize(n + 1);
    for (int u = 1; u <= n; u++) cin >> color[u];
    g.resize(n + 1);
    for (int i = 0, u, v; i < n - 1; i++) {
        cin >> u >> v;
        g[u].push_back(v); g[v].push_back(u);
    }
    sz.resize(n + 1); heavy.resize(n + 1);
    cnt.resize(n + 1, 0); ans.resize(n + 1);

    dfs_sz(1, 0);             // root at node 1
    dfs_dsu(1, 0, true);
    for (int u = 1; u <= n; u++) cout << ans[u] << ' ';
    cout << '\n';
}
\end{lstlisting}

The colour-frequency state contains \texttt{mxFreq} and \texttt{sumOfColors}. When a colour reaches a new maximum, it is the only colour with that frequency. When it reaches the current maximum, add it to the sum. During a complete clear, the counts are removed and these aggregate values reset to zero.

\section{Generic DSU on Tree template}

The problem-specific parts are the state, what happens when a node is added or removed, and the answer recorded at each node. The tree normally must be rooted first; \texttt{dfs\_sz} finds subtree sizes and the heavy child, then \texttt{dfs\_dsu} performs the algorithm.

\begin{lstlisting}
int n;
vector<vector<int>> g;
vector<int> color, sz, heavy, cnt;
vector<long long> ans;

void dfs_sz(int u, int p) {
    sz[u] = 1; heavy[u] = -1;
    int largest = 0;
    for (int v : g[u]) if (v != p) {
        dfs_sz(v, u);
        sz[u] += sz[v];
        if (sz[v] > largest) largest = sz[v], heavy[u] = v;
    }
}

void add(int u, int p) {
    // Update all problem-specific state for node u here.
    // Example for colour frequencies:
    cnt[color[u]]++;
    for (int v : g[u]) if (v != p) add(v, u);
}

void remove(int u, int p) {
    // This must exactly reverse add().
    cnt[color[u]]--;
    for (int v : g[u]) if (v != p) remove(v, u);
}

void clearState() {
    // Reset extra aggregate state here, e.g. distinct, maximumFrequency, sum.
}

void dfs_dsu(int u, int p, bool keep) {
    for (int v : g[u])
        if (v != p && v != heavy[u]) dfs_dsu(v, u, false);
    if (heavy[u] != -1) dfs_dsu(heavy[u], u, true);
    for (int v : g[u])
        if (v != p && v != heavy[u]) add(v, u);

    // Add only u itself, then calculate the problem-specific answer.
    cnt[color[u]]++;
    ans[u] = 0; // Replace with the current statistic.

    if (!keep) {
        remove(u, p);
        clearState();
    }
}
\end{lstlisting}

\section{Mental model}

For a node \texttt{u} with light children and one heavy child:
\begin{enumerate}
\item Solve every light subtree and throw its data away.
\item Solve the heavy subtree with \texttt{keep = true}.
\item Add the light subtrees back to the retained heavy data.
\item Add \texttt{u}, answer the query for \texttt{u}, and clear only if the parent does not need the data.
\end{enumerate}
The core line is \texttt{dfs\_dsu(heavy[u], u, true)}: keeping the largest subtree means any node is added and removed only logarithmically many times.

\section{Common customizations}

\subsection{Number of distinct colours}
\begin{lstlisting}
int distinct = 0;
// Add:
if (cnt[color[u]] == 0) distinct++;
cnt[color[u]]++;
// Remove:
cnt[color[u]]--;
if (cnt[color[u]] == 0) distinct--;
// Answer:
ans[u] = distinct;
\end{lstlisting}

\subsection{Maximum frequency}
\begin{lstlisting}
int mxFreq = 0;
// Add:
cnt[color[u]]++;
mxFreq = max(mxFreq, cnt[color[u]]);
// Answer:
ans[u] = mxFreq;
\end{lstlisting}
If colours are removed individually, updating \texttt{mxFreq} with only \texttt{max} is insufficient. Use a frequency-of-frequencies structure, or reset it when the full state is cleared as appropriate for the implementation.

\subsection{Sum of colours with maximum frequency}
\begin{lstlisting}
int mxFreq = 0;
long long sum = 0;
// Add colour c:
cnt[c]++;
if (cnt[c] > mxFreq) mxFreq = cnt[c], sum = c;
else if (cnt[c] == mxFreq) sum += c;
// Answer:
ans[u] = sum;
\end{lstlisting}
For fully dynamic individual removals, a frequency-of-frequencies structure makes the maximum and sum safe to maintain.

\section{When to use it}
\begin{itemize}
\item \textbf{DFS + brute force}: small constraints.
\item \textbf{Small-to-large map merging}: subtree data that naturally belongs in maps or sets.
\item \textbf{DSU on Tree}: many subtree statistics with efficient add/remove operations.
\item \textbf{Euler Tour + Fenwick/Segment Tree}: when a subtree can be treated as a contiguous array range.
\item \textbf{Mo's on Tree}: more complicated offline tree queries.
\item \textbf{Heavy-Light Decomposition}: path queries and updates.
\end{itemize}

Rule of thumb: if a problem asks for a statistic over every node's subtree, and the statistic supports efficient \texttt{add(node)} and \texttt{remove(node)}, DSU on Tree is a strong candidate.
